\documentclass{article}
\usepackage[utf8]{inputenc}
\usepackage[T1]{fontenc}
\usepackage{lmodern}
\usepackage[a4paper, margin=1in]{geometry}
\usepackage{minted}
\usepackage[spanish]{babel}
\usepackage{enumitem}

% Configuración global para el resaltado de código R
\setminted[r]{
    frame=lines,
    linenos,
    fontsize=\large,
    xleftmargin=2em
}

\large
\title{Taller Práctico de R}

\begin{document}

\begin{titlepage}
    \begin{center}
        \line(1,0){300}\\[0.65cm]
        \huge{\bfseries Asignación de variables y operaciones aritméticas}\\[0.3cm]
        \line(1,0){300}\\[0.5cm]
        \textsc{\Large Bioestadística Intermedia}\\[0.5cm]
        \textsc{\LARGE \today}\\[4.5cm]     
    \end{center}
    
    \begin{flushright}
        \textsc{\Large Dr. Ricardo Iván Bravo Chávez}\\[0.3cm]
        \textsc{\Large Área Académica de Gerontología}\\[0.3cm]
        \textsc{\Large Maestría en Ciencias Biomédicas}\\[0.5cm]
    \end{flushright}
\end{titlepage}

\section*{Ejercicio 1}

\subsection*{Enunciado}
\textquestiondown Cuál de las siguientes expresiones en R calcula 
correctamente la siguiente operación matemática 
\begin{center}
    \large
$\frac{10+5}{2\times3}$
\end{center}

%\subsection*{Demostración en R}
\begin{minted}{r}
# Elija la opción correcta
10 + 5 / 2 * 3      # Opción a)
(10 + 5) / 2 * 3    # Opcion b)
(10 + 5) / (2 * 3)  # Opción c)
10 + (5 / 2) * 3    # Opción d)
\end{minted}

\section*{Ejercicio 2}
Un estudiante intentó calcular $5 \times 4$ escribiendo en la consola: 5(4). 
R devolvió el siguiente error 

\begin{minted}{r}
# Piense en el error
5(4)
"Error in 5(4): attempt to apply non-function"
\end{minted}

\textquestiondown Por qué ocurre este error? \textquestiondown Cómo 
debe escribirse la expresión correctamente en R?

% \subsection*{Solución en R}
% \begin{minted}{r}
% # Solución 2
% v <- c(1, 4, 9, 16, 25)
% \end{minted}

\section*{Ejercicio 3}
Analiza la siguiente secuencia de código en R y determina cuál será 
el valor final de la variable

\begin{minted}{r}
# Determina el valor de la variable
a <- 10
b <- a * 2
a <- 5
\end{minted}

\begin{enumerate}[label=\alph*]
    \item 10
    \item 20
    \item 5
    \item R lanza una advertencia
\end{enumerate}

\section*{Ejercicio 4}
Imagina que creaste las siguientes variables para probar un cálculo: 

\begin{minted}{r}
temp_1 <- 45
temp_2 <- 90
resultado <- temp_1 + temp_2
\end{minted}

\subsection*{4.1}
Escribe la instrucción en R para borrar únicamente las variables 
temporales (\textbf{temp\_1} y \textbf{temp\_2}) sin eliminar el resultado. 

\subsection*{4.2}
Escribe un comando que te permite listar todos los objetos que están 
actualmente guardados en tu entorno de trabajo (\textit{Enviroment}).

\section*{Ejercicio 5}
Tienes la siguiente serie de calificaciones de un estudiante: 85, 90, 78, 92, 88, 60.
\begin{itemize}
    \item Escribe el código en R para almacenar estas calificaciones en un vector llamado "notas".
    \item Calcula la nota \textit{promedio} utilizando las funciones incorporadas \textbf{sum()} y \textbf{length()}
    \item Verifica tu resultado anterior utilizando la función \textbf{mean()}.
\end{itemize}

\section*{Ejercicio 6}
Un comerciante tiene los precios de 4 productos almacenados en un vector:
\begin{minted}{r}
    precios <- c(100,250,80,150)
\end{minted}

Debido a la inflación, necesita aplicar un incremento del \textbf{15\%} a todos 
los productos y además sumar un costo fijo de \textbf{\$5} a cada uno. Escribe 
una sola línea de código en R para calcular el nuevo precio total 
de cada producto y guardarlo en un vector llamado \textbf{nuevos\_precios}.

\section*{Ejercicio 7}
Dado el vector
\begin{minted}{r}
dias <- c("Lunes", "Martes", "Miércoles", "Jueves", "Viernes", 
"Sábado", "Domingo")
\end{minted}
\textquestiondown Qué instrucción utilizarías para extraer únicamente 
los días del fin de semana (``Sabado'' y ``Domingo'')
\textquestiondown Qué devuelve la instrucción \textbf{dias[-1]}?



\end{document}
